% ── Chapter 3: Methodology ────────────────────────────────────
\chapter{ระเบียบวิธีวิจัย (Methodology)}

\section{ภาพรวมสถาปัตยกรรมระบบ (System Architecture Overview)}

ระบบที่นำเสนอพัฒนาในรูปแบบเว็บแอปพลิเคชันครบวงจร (Full-stack Web Application) ตามสถาปัตยกรรม Client-server โดยแบ่งความรับผิดชอบออกเป็น 3 ส่วนหลักอย่างชัดเจน ได้แก่ ส่วนติดต่อผู้ใช้งาน (User Interface), AI Inference Pipeline และชั้นจัดการข้อมูล (Data Persistence Layer) การแบ่งส่วนนี้ทำให้พัฒนาและทดสอบแต่ละส่วนได้อย่างอิสระ และยังให้ Backend ซึ่งรับภาระการประมวลผล AI ปรับขยายหรืออัปเกรดได้โดยไม่กระทบฝั่ง Frontend

สำหรับฝั่ง Frontend ได้รับการพัฒนาด้วย React ซึ่งเป็นเฟรมเวิร์ก JavaScript แบบ Component-based ที่ทำงานประสานกับ HTML \cite{sharma2018html,mozilla2025html} ทำหน้าที่เป็นหน้าจอหลัก (Interface) ให้ผู้ใช้สามารถอัปโหลดรูปภาพ หรือสั่งถ่ายภาพจากกล้องได้แบบเรียลไทม์ พร้อมแสดงผลลัพธ์ผ่านหน้าแดชบอร์ดแบบโต้ตอบ (Interactive Dashboard) ในขณะที่ฝั่ง Backend ถูกพัฒนาด้วยภาษา Python \cite{siva2023python} ร่วมกับเฟรมเวิร์ก FastAPI \cite{tioman2024fastapi} ซึ่งมีประสิทธิภาพสูงในการประมวลผลแบบอซิงโครนัส (Asynchronous) ทำหน้าที่ให้บริการ RESTful API แก่ฝั่ง Frontend 
เมื่อได้รับรูปภาพผ่าน API ระบบ Backend จะเรียกใช้ AI Pipeline ทันที และส่งผลลัพธ์กลับในรูปแบบ JSON ประกอบด้วย พิกัดการตรวจจับวัตถุทั้งหมด, ค่าเส้นผ่านศูนย์กลาง, ชื่อยาปฏิชีวนะจาก OCR และผลจำแนกความไวต่อยา (S/I/R) นอกจากนี้ ระบบยังผสานการทำงานกับฐานข้อมูลเชิงสัมพันธ์ SQLite \cite{sqlite} เพื่อเก็บบันทึกประวัติการวิเคราะห์, ข้อมูลเซสชันการใช้งาน และจัดเก็บตารางเกณฑ์จุดตัด (Breakpoints) จากมาตรฐาน CLSI 2025 และ EUCAST 2023 เพื่อให้ระบบสามารถดึงข้อมูลมาเปรียบเทียบและแปลผลได้อย่างรวดเร็ว

ภาพรวมของกระบวนการวิเคราะห์ภาพ (Pipeline) ประกอบด้วย 5 ขั้นตอนหลัก ดังนี้:
\begin{enumerate}
  \item \textbf{การนำเข้าและประมวลผลภาพเบื้องต้น (Image Acquisition and Preprocessing):} ภาพถ่ายที่ผู้ใช้นำเข้าจะถูกตรวจสอบความถูกต้อง ปรับสัดส่วนให้ตรงกับความละเอียดที่โมเดล AI ต้องการ และผ่านกระบวนการปรับแต่งภาพเบื้องต้น (เช่น ปรับแสง สี) ให้สอดคล้องกับเงื่อนไขของโมเดล
  \item \textbf{การตรวจจับแผ่นยาและระบุชื่อด้วย OCR (Antibiotic Disk Detection and OCR Labeling):} โมเดล Object Detection จะระบุตำแหน่งของแผ่นยาปฏิชีวนะแต่ละแผ่นในภาพ จากนั้นภาพแผ่นยาจะถูกตัด (Crop) แยกออกมาทีละแผ่น และส่งให้ระบบ OCR อ่านรหัสตัวอักษรและตัวเลขเพื่อระบุชนิดของยา
  \item \textbf{การตรวจจับและวัดขนาด ZOI (ZOI Detection and Diameter Measurement):} โมเดลจะตรวจหา ZOI และวาดกรอบ (Bounding Box) ล้อมรอบ จากนั้นขนาดความกว้างและความสูงของกรอบจะถูกนำมาคำนวณเป็นค่าประมาณเส้นผ่านศูนย์กลางในหน่วยพิกเซล
  \item \textbf{การปรับเทียบสัดส่วนภาพ (Pixel-to-Millimeter Calibration):} ระบบจะแปลงค่าเส้นผ่านศูนย์กลางจากหน่วยพิกเซลให้เป็นหน่วยมิลลิเมตรทางกายภาพ โดยอาศัยเส้นผ่านศูนย์กลางมาตรฐาน 6 มม. ของแผ่นยาปฏิชีวนะที่ปรากฏในภาพเป็นตัวเลขอ้างอิง (Reference)
  \item \textbf{การเปรียบเทียบเกณฑ์และการแปลผล (Breakpoint Lookup and Susceptibility Classification):} ขนาดเส้นผ่านศูนย์กลาง ZOI ที่แปลงเป็นมิลลิเมตรแล้ว จะถูกนำไปเทียบกับฐานข้อมูลตาราง Breakpoint ของ CLSI หรือ EUCAST (ตามที่ผู้ใช้เลือก) โดยอ้างอิงจากชื่อยาและชนิดของเชื้อแบคทีเรีย เพื่อสรุปผลการจำแนกความไวต่อยาออกเป็น 3 ระดับ คือ ไวต่อยา (S), กึ่งไว (I) หรือ ดื้อยา (R)
\end{enumerate}

\begin{figure}[H]
  \centering
  \safeimg[width=\textwidth]{figures/pipeline_flowchart.png}{End-to-end analysis pipeline}
  \caption{Pipeline การวิเคราะห์รูปภาพแบบ End-to-end ของระบบ AST อัตโนมัติ ตั้งแต่การนำเข้าภาพผ่านการตรวจจับ, OCR, การปรับเทียบ และการจำแนกความไวต่อยาตาม Breakpoint}
  \label{fig:pipeline}
\end{figure}

\begin{figure}[H]
  \centering
  \safeimg[width=0.92\textwidth]{figures/system_architecture.png}{System Architecture Diagram}
  \caption{สถาปัตยกรรมระบบโดยรวมของเว็บแอปพลิเคชัน AST อัตโนมัติ แสดงให้เห็นส่วน React frontend, FastAPI backend, AI inference pipeline และชั้นฐานข้อมูล SQLite}
  \label{fig:sys_arch}
\end{figure}

\begin{figure}[H]
  \centering
  \safeimg[width=0.82\textwidth]{figures/_review/use_case_diagram.png}{Use case diagram}
  \caption{แผนภาพ Use Case แสดงการโต้ตอบระหว่างผู้ใช้งาน (บุคลากรทางการแพทย์)
    กับระบบวิเคราะห์ AST อัตโนมัติ ครอบคลุมการจัดการบัญชี การวิเคราะห์ภาพ และการดูผลลัพธ์}
  \label{fig:use_case}
\end{figure}

\section{การรวบรวมและการเตรียมชุดข้อมูล (Dataset Collection and Preparation)}

การจัดเตรียมชุดข้อมูลฝึกสอน (Training Dataset) ที่มีทั้งคุณภาพและความหลากหลายสูง ถือเป็นหัวใจสำคัญที่จะทำให้โมเดล Deep Learning สามารถเรียนรู้และนำไปใช้งานจริง (Generalization) ได้อย่างมีประสิทธิภาพ อย่างไรก็ดี การรวบรวมข้อมูลภาพถ่ายจานเพาะเชื้อ AST ทางคลินิกมีความท้าทายอย่างมาก เนื่องจากภาพถ่ายจริงมักถูกจัดเป็นข้อมูลส่วนบุคคลหรือกรรมสิทธิ์ของสถานพยาบาล ทำให้มีชุดข้อมูลภาพ AST ที่ได้รับการ Label (Annotated) เผยแพร่สู่สาธารณะน้อยมาก โครงงานนี้จึงแก้ปัญหาด้วยกลยุทธ์การผสานรวมชุดข้อมูลจาก 3 แหล่ง เพื่ออุดช่องโหว่ซึ่งกันและกัน ได้แก่ 1) ภาพถ่ายจริงจากห้องปฏิบัติการ 2) ภาพถ่ายจริงที่ผ่านกระบวนการทำ Data Augmentation และ 3) ภาพสังเคราะห์ (Synthetic Images) ที่สร้างขึ้นด้วยโปรแกรมคอมพิวเตอร์

\subsection{การรวบรวมภาพถ่ายจริง (Real Image Collection)}

ภาพถ่ายจานเพาะเชื้อ AST จริงทั้งหมด 287 ภาพ ถูกรวบรวมมาจาก 3 แหล่งหลัก แหล่งแรกคือฐานข้อมูลสาธารณะ DRYAD \cite{giske2024dryad} ซึ่งสนับสนุนแนวคิดการแบ่งปันข้อมูลวิจัยแบบเปิดกว้าง \cite{iom1996} โดยได้ภาพจาน Disk Diffusion ที่มีป้ายกำกับแล้วจำนวน 225 ภาพ ถ่ายจากสภาพแวดล้อมห้องปฏิบัติการที่แตกต่างกัน แหล่งที่สองคือภาพชุดข้อมูลส่วนตัวจากงานวิจัย CDS (Clinical and Demographic Study) จำนวน 50 ภาพ และแหล่งสุดท้ายคือภาพจากวิทยาลัยแพทยศาสตร์ศรีสวางควัฒน ราชวิทยาลัยจุฬาภรณ์ (CRA) จำนวน 32 ภาพ

ภาพถ่ายทั้งหมดนี้ครอบคลุมความหลากหลายของสายพันธุ์แบคทีเรีย จำนวนแผ่นยาปฏิชีวนะต่อจาน (ตั้งแต่ 1 ถึง 10 แผ่น) และระดับการซ้อนทับกันของ ZOI ซึ่งสะท้อนถึงสภาพความเป็นจริงทางคลินิกที่ระบบต้องเผชิญ การรวบรวมข้อมูลจากหลายแหล่งช่วยเพิ่มความแปรผันให้กับชุดข้อมูล ทั้งในแง่ของสภาพแสง รุ่นกล้องที่ใช้ถ่าย ยี่ห้อของวุ้น (Agar) และขนาดของจานเพาะเชื้อ ซึ่งล้วนเป็นผลดีต่อการเพิ่มความสามารถในการสรุปนัย (Generalizability) ของโมเดล

\subsection{การเพิ่มขยายข้อมูล (Data Augmentation)}

เพื่อขยายขนาดชุดข้อมูลฝึกสอน (Training Set) จากภาพต้นฉบับ 287 ภาพ โครงงานได้ประยุกต์ใช้เทคนิคการเพิ่มข้อมูล (Data Augmentation) ทั้งในเชิงเรขาคณิตและเชิงแสงผ่านแพลตฟอร์ม Roboflow \cite{roboflow} กระบวนการนี้ครอบคลุมตั้งแต่ การพลิกภาพแนวตั้งและแนวนอน (Horizontal/Vertical Flipping), การหมุนภาพแบบสุ่ม ($\pm 15^\circ$), การปรับความสว่าง ($\pm 25\%$), การปรับความเปรียบต่าง ($\pm 20\%$) และการเติมสัญญาณรบกวน (Gaussian Noise โดยมีค่าเบี่ยงเบนมาตรฐานไม่เกิน 5\% ของความเข้มพิกเซลสูงสุด) ระบบจะสุ่มเลือกเทคนิคเหล่านี้มาใช้สร้างภาพใหม่ โดยพิกัดของ Bounding Box จะถูกคำนวณและปรับเปลี่ยนตามรูปทรงของภาพที่เปลี่ยนไปโดยอัตโนมัติ ทำให้ได้ชุดข้อมูลสำหรับฝึกสอนเพิ่มขึ้นเป็น 543 ภาพ (เพิ่มขึ้นเกือบเท่าตัวจากภาพจริง)

ข้อควรระวังประการสำคัญคือ เทคนิค Data Augmentation จะถูกจำกัดการใช้งานเฉพาะกับข้อมูลกลุ่มที่ใช้สอนโมเดล (Training Partition) เท่านั้น ส่วนข้อมูลกลุ่มตรวจสอบ (Validation Partition) และกลุ่มทดสอบ (Test Partition) จะใช้ภาพจริงดั้งเดิมที่ไม่มีการดัดแปลงใดๆ เพื่อให้มั่นใจว่าค่าชี้วัดประสิทธิภาพที่ได้จะสะท้อนถึงความสามารถของโมเดลเมื่อต้องทำงานกับภาพจริงในห้องปฏิบัติการ นอกจากนี้ รูปแบบการทำ Augmentation จะถูกจำกัดให้อยู่ในขอบเขตที่สมจริงและสามารถเกิดขึ้นได้จริงในห้องปฏิบัติการ โดยหลีกเลี่ยงการบิดเบือนมุมมองที่รุนแรงเกินไป (Extreme Perspective Distortion) หรือการสลับค่าสี (Color Inversion) ซึ่งจะทำให้ภาพสูญเสียความสมจริงและส่งผลเสียต่อการเรียนรู้ของโมเดล

\subsection{การสร้างข้อมูลสังเคราะห์ (Synthetic Data Generation)}

เพื่อแก้ปัญหาการขาดแคลนภาพถ่ายจริงในกรณีศึกษาที่มี ZOI ซ้อนทับกันอย่างซับซ้อน หรือกรณีที่ ZOI มีรูปทรงผิดปกติ โครงงานจึงได้ใช้โปรแกรมสร้างภาพจานเพาะเชื้อสังเคราะห์ (Synthetic Images) ขึ้นมาเพิ่มเติมอีก 2,000 ภาพ ผ่านกระบวนการ Procedural Pipeline ที่พัฒนาขึ้นด้วยภาษา Python (ใช้ไลบรารี Pillow และ OpenCV) ภาพสังเคราะห์เหล่านี้ถูกออกแบบมาให้มีลักษณะทางสายตาคล้ายคลึงกับจานวุ้น Mueller-Hinton ของจริงมากที่สุด โดยมีองค์ประกอบครบถ้วน ได้แก่ ลวดลายพื้นผิวของวุ้น (Texture) ที่เกิดจากการผสมผสานลวดลายตัวอย่าง แผ่นยาปฏิชีวนะทรงกลมในตำแหน่งที่กำหนด และ ZOI ที่ถูกเรนเดอร์ให้มีขอบฟุ้งเบลอ (Gaussian Blur) เพื่อจำลองรูปแบบการเจริญเติบโตของแบคทีเรียที่ค่อยๆ จางหายไปบริเวณขอบ ZOI ดังที่ปรากฏในภาพถ่ายจริง

\begin{figure}[H]
  \centering
  \safeimg[width=0.58\textwidth]{figures/_review/synthetic_plate_example.png}{Synthetic plate example}
  \caption{ตัวอย่างภาพ Synthetic ที่สร้างขึ้นด้วยโปรแกรม แสดงจานเพาะเชื้อจำลองพร้อมแผ่นยาปฏิชีวนะ
    และ ZOI ที่มีขนาดหลากหลาย ใช้สำหรับเพิ่มปริมาณและความหลากหลายของชุดข้อมูลฝึกฝน}
  \label{fig:synthetic_example}
\end{figure}

รูปภาพสังเคราะห์ 2,000 ภาพถูกแบ่งออกเป็นสองประเภทในจำนวนที่เท่ากัน ประเภทแรกจำนวน 1,000 ภาพใช้ตำแหน่งแผ่นยาปฏิชีวนะแบบคงที่ตามรูปแบบจานเพาะเชื้อที่แนะนำโดย CLSI สำหรับการวางแผ่นยา 6 แผ่นและ 8 แผ่น เพื่อจำลองรูปแบบมาตรฐานที่เป็นระบบในขั้นตอนการทำงานของห้องปฏิบัติการทั่วไป ประเภทที่สองจำนวน 1,000 ภาพใช้การวางแผ่นยาแบบสุ่มภายในขอบเขตของจาน เพื่อจำลองการจัดวางที่เป็นระบบน้อยกว่าซึ่งพบได้บ่อยในห้องปฏิบัติการวิจัยและการเรียนการสอน ทั้งสองประเภทมีการสุ่มขนาดเส้นผ่านศูนย์กลางของบริเวณยับยั้งจากช่วงค่าทางคลินิกที่สมจริง (ตั้งแต่ประมาณ 10\,มม. ถึง 40\,มม.) โดยจงใจเพิ่มสัดส่วนของสถานการณ์ที่บริเวณยับยั้งมีการซ้อนทับกัน (ซึ่งบริเวณยับยั้งของแผ่นยาที่อยู่ใกล้กันบางส่วนตัดกัน) เพื่อให้แน่ใจว่าโมเดลได้เรียนรู้กรณีที่ท้าทายแต่มีความสำคัญทางคลินิกนี้ในระหว่างการฝึกฝน ลวดลายพื้นหลัง, การจำลองการไล่ระดับแสง และรูปแบบการเจริญเติบโตของโคโลนีสังเคราะห์ก็ถูกทำให้หลากหลายเพื่อเพิ่มความหลากหลายทางสายตา

\subsection{การแบ่งชุดข้อมูลขั้นสุดท้าย (Final Dataset Split)}

\begin{table}[H]
  \centering
  \caption{องค์ประกอบของชุดข้อมูลและการแบ่งส่วนสำหรับการฝึกฝน/การตรวจสอบ/การทดสอบ ที่ใช้ในการฝึกฝนและประเมินผลโมเดล}
  \label{tab:dataset_split}
  \begin{tabular}{lcccc}
    \toprule
    \textbf{แหล่งข้อมูล (Source)} & \textbf{ทั้งหมด} & \textbf{ฝึกฝน (Train)} & \textbf{ตรวจสอบ (Validation)} & \textbf{ทดสอบ (Test)} \\
    \midrule
    DRYAD                              & 225   & ---   & ---   & ---   \\
    CDS                                & 50    & ---   & ---   & ---   \\
    CRA (ราชวิทยาลัยจุฬาภรณ์)           & 32    & ---   & ---   & ---   \\
    \textit{รวมรูปภาพจริง (ต้นฉบับ)}     & 287   & ---   & ---   & ---   \\
    \textit{รวมรูปภาพจริง (เพิ่มข้อมูล)}  & 543   & 433   & 55    & 55    \\
    ข้อมูลสังเคราะห์ (ตำแหน่งคงที่)       & 1,000 & \multicolumn{3}{c}{ใช้ในการฝึกฝนเท่านั้น} \\
    ข้อมูลสังเคราะห์ (ตำแหน่งไม่คงที่)     & 1,000 & \multicolumn{3}{c}{ใช้ในการฝึกฝนเท่านั้น} \\
    \bottomrule
  \end{tabular}
\end{table}

รูปภาพสังเคราะห์ถูกนำมาใช้เพื่อการฝึกฝนเท่านั้น เนื่องจากถูกสร้างขึ้นมาเพื่อเสริมการเรียนรู้ของโมเดลในรูปแบบที่หลากหลาย มากกว่าที่จะใช้ประเมินประสิทธิภาพการทำงานทั่วไป ตัวชี้วัดประสิทธิภาพทั้งหมดที่รายงานในบทที่ 4 คำนวณจากชุดทดสอบรูปภาพจริง (55 ภาพ) เพื่อให้แน่ใจว่าตัวเลขประสิทธิภาพที่รายงานสะท้อนถึงพฤติกรรมของโมเดลกับภาพถ่ายในห้องปฏิบัติการที่เป็นของจริงและไม่ได้ถูกตกแต่ง

\section{การระบุป้ายกำกับ (Annotation)}

ป้ายกำกับ Ground-truth ถูกสร้างขึ้นในรูปแบบ Pascal VOC XML โดยระบุพิกัด Bounding box สำหรับแต่ละบริเวณยับยั้งเชื้อและแต่ละแผ่นยาปฏิชีวนะที่ปรากฏในรูปภาพ รูปแบบการระบุป้ายกำกับนี้จะเข้ารหัสค่าสี่ค่าต่อวัตถุ ได้แก่ พิกัดพิกเซลของมุมบนซ้ายและมุมล่างขวาของ Bounding box พร้อมกับป้ายกำกับคลาส (\texttt{antibiotic} สำหรับแผ่นยา และ \texttt{disk-zone} สำหรับบริเวณยับยั้งเชื้อ)

การระบุป้ายกำกับเบื้องต้นทำโดยทีมงานโปรเจกต์โดยใช้ LabelImg \cite{labelimg} ซึ่งเป็นเครื่องมือระบุป้ายกำกับรูปภาพแบบกราฟิกที่เป็นโอเพนซอร์ส เพื่อให้แน่ใจว่ามีความถูกต้องทางคลินิก การวาง Bounding box ทั้งหมดสำหรับคลาสบริเวณยับยั้งเชื้อได้รับการตรวจสอบและแก้ไขโดยนักจุลชีววิทยาที่ผ่านการฝึกฝนมาแล้ว ซึ่งเป็นผู้ตรวจสอบป้ายกำกับแต่ละรายการเปรียบเทียบกับการวัดจานจริงโดยใช้ Calipers ขั้นตอนการตรวจสอบสองชั้นนี้เป็นสิ่งจำเป็นสำหรับการสร้าง Ground-truth ที่เชื่อถือได้ เนื่องจากความหมายของขอบเขตบริเวณยับยั้งต้องใช้ความเชี่ยวชาญเฉพาะทาง: ขอบที่ถูกต้องของบริเวณยับยั้งถูกกำหนดให้เป็นจุดที่การยับยั้งการเจริญเติบโตแบบ Confluent growth สิ้นสุดลงโดยสมบูรณ์ (ไม่ใช่บริเวณที่มีการเกิด Micro-colony ซึ่ง CLSI แนะนำให้เจ้าหน้าที่ละเว้น) จากนั้นป้ายกำกับจะถูกนำเข้าสู่ Roboflow เพื่อจัดการชุดข้อมูล, กำหนด Pipeline การเพิ่มข้อมูล และส่งออกในรูปแบบที่เฟรมเวิร์กการฝึกฝนของแต่ละโมเดลต้องการ (Pascal VOC XML สำหรับ Faster R-CNN และ YOLO TXT สำหรับ YOLOv11n)

\begin{figure}[H]
  \centering
  \safeimg[width=0.82\textwidth]{figures/_review/annotation_roboflow.png}{Roboflow annotation screenshot}
  \caption{หน้าจอการระบุป้ายกำกับ (Annotation) บนแพลตฟอร์ม Roboflow แสดง Bounding box
    สำหรับคลาส \texttt{antibiotic} (ดิสก์ยาปฏิชีวนะ สีแดง) และ \texttt{disk-zone}
    (ZOI สีม่วง) พร้อมสถิติจำนวนตัวอย่างในแต่ละคลาส}
  \label{fig:annotation}
\end{figure}

\section{การฝึกฝนโมเดล (Model Training)}

\subsection{Faster R-CNN}

Faster R-CNN ถูกพัฒนาโดยใช้ไลบรารี \texttt{torchvision} ของ PyTorch พร้อมกับ Backbone แบบ ResNet-50 ที่ผ่านการฝึกฝนเบื้องต้น (Pretrained) บนชุดข้อมูล ImageNet \cite{he2016deep} การใช้ Backbone ที่ผ่านการฝึกฝนมาแล้วเป็นสิ่งสำคัญมากสำหรับประสิทธิภาพของชุดข้อมูลภาพทางการแพทย์ที่มีขนาดจำกัดเช่นนี้: การฝึกฝนเบื้องต้นบน ImageNet ช่วยให้ Backbone มีตัวแทนคุณลักษณะ (Feature representations) ที่แยกแยะโครงสร้างรูปภาพระดับต่ำ (ขอบ, พื้นผิว, การไล่ระดับสี) ซึ่งสามารถถ่ายโอนมายังโดเมน AST ได้อย่างมีประสิทธิภาพ ช่วยเร่งการลู่เข้าหาคำตอบ (Convergence) และปรับปรุงความแม่นยำขั้นสุดท้ายได้ดีกว่าการเริ่มฝึกฝนจากการสุ่มค่าเริ่มต้น

ส่วนหัวของการตรวจจับ (Detection head ทั้ง RPN และ Classification/regression head) ถูกกำหนดค่าเริ่มต้นใหม่และฝึกฝนตั้งแต่ต้น การฝึกฝนใช้อัลกอริทึมการเพิ่มประสิทธิภาพ Stochastic Gradient Descent (SGD) โดยมี Learning rate เท่ากับ $0.005$, Momentum เท่ากับ $0.9$ และ Weight decay เท่ากับ $5 \times 10^{-4}$ (คือ \texttt{0.0005}) เพื่อรักษาสมดุลของโมเดลและป้องกันการ Overfitting กับชุดข้อมูลรูปภาพจริงที่มีอยู่อย่างจำกัด ค่า Learning rate จะถูกปรับลดลงโดยใช้ \texttt{StepLR} scheduler โดยมี \texttt{step\_size=3} และ \texttt{gamma=0.1} ซึ่งจะลด Learning rate ลง 10 เท่าทุกๆ 3 Epochs รูปภาพทั้งหมดถูกปรับขนาดเป็น $640 \times 640$ พิกเซลก่อนนำเข้าโมเดล และใช้ Batch size เท่ากับ 8

ก่อนที่จะเริ่มการฝึกฝนแบบเต็มรูปแบบ ได้มีการรันการทดลองฝึกฝนเบื้องต้นบนชุดย่อยของชุดข้อมูลที่ขยายขนาดขึ้นเรื่อยๆ เพื่อศึกษาพฤติกรรมการเรียนรู้ของโมเดลและวินิจฉัยสภาวะ Underfitting และ Overfitting การทดลองเหล่านี้ช่วยในการเลือก Hyperparameter ขั้นสุดท้าย โดยเฉพาะจำนวน Epochs ในการฝึกฝน (10 Epochs) และตารางเวลาของ Learning rate โดยแสดงให้เห็นว่า Epoch ใดที่ Loss ของการตรวจสอบ (Validation loss) มีค่าต่ำสุดก่อนที่จะเริ่มสูงขึ้นอีกครั้ง การฝึกฝนขั้นสุดท้ายใช้เวลาประมาณ 174 นาทีบนฮาร์ดแวร์ที่มีอยู่ ซึ่งสะท้อนถึงต้นทุนการคำนวณของสถาปัตยกรรม Inference แบบสองขั้นตอน

\subsection{YOLOv11n}

YOLOv11n ถูกฝึกฝนโดยใช้เฟรมเวิร์กการฝึกฝนของ Ultralytics ซึ่งให้บริการ Python API ระดับสูงสำหรับการกำหนดค่าเริ่มต้นโมเดล, การดำเนินการลูปการฝึกฝน และการบันทึก Checkpoint \cite{wang2024yolov11} ตัวเลือกย่อยรุ่น Nano (n) ถูกเลือกเพื่อเน้นความเร็วในการประมวลผล (Inference speed) และความสามารถในการนำไปใช้งานบนระบบที่ไม่มี GPU แยก ในขณะที่ยังคงรักษาสมรรถนะของโมเดลเพียงพอที่จะเรียนรู้คุณลักษณะทางสายตาที่ค่อนข้างจำกัดของรูปภาพจานเพาะเชื้อ AST (คลาสวัตถุสองคลาสที่มีรูปร่างที่ชัดเจน: แผ่นยาทรงกลม และบริเวณยับยั้งทรงกลม)

น้ำหนักของโมเดลถูกกำหนดค่าเริ่มต้นจาก Checkpoint ที่ผ่านการฝึกฝนบน COCO ของ Ultralytics สำหรับ YOLOv11n การเรียนรู้แบบถ่ายโอน (Transfer learning) จาก COCO ซึ่งเป็นฐานข้อมูลการตรวจจับวัตถุในธรรมชาติขนาดใหญ่ที่มีวัตถุ 80 ประเภท ช่วยให้โมเดลมีความสามารถในการตรวจจับวัตถุทั่วไป (การดึงคุณลักษณะ, การระงับค่าที่ไม่ใช่ค่าสูงสุดหรือ Non-maximum suppression, การพยากรณ์แบบ Anchor-free) ซึ่งจะถูกปรับจูน (Fine-tuning) ให้เข้ากับโดเมนเฉพาะอย่าง AST การฝึกฝนแบบ Fine-tuning ถูกรันเป็นจำนวน \textbf{200 Epochs} บนการแบ่งชุดข้อมูลเดียวกันกับ Faster R-CNN โดยรูปภาพถูกปรับขนาดเป็น $640 \times 640$ พิกเซล และใช้ Batch size เท่ากับ 8 โดยได้ปิดฟังก์ชันการเพิ่มข้อมูลในตัวของ Ultralytics (\texttt{augment=False}) เพื่อให้โมเดลได้รับการฝึกฝนด้วยรูปภาพที่สะอาดและไม่มีการเพิ่มข้อมูล ซึ่งสอดคล้องกับกลยุทธ์การเพิ่มข้อมูลที่ใช้ในขั้นตอนการเตรียมชุดข้อมูล การฝึกฝนด้วยชุดข้อมูลเต็มรูปแบบในการกำหนดค่านี้ใช้เวลาประมาณ 200 นาที

YOLOv11n ตรวจจับวัตถุสองคลาส ได้แก่ \texttt{antibiotic} (แผ่นยาที่มีตัวยา ซึ่งทำหน้าที่เป็นเป้าหมายของ OCR และเป็นตัวอ้างอิงในการปรับเทียบในภาพ) และ \texttt{disk-zone} (บริเวณยับยั้งโดยรอบ ซึ่งใช้สำหรับวัดเส้นผ่านศูนย์กลาง ZOI) ในระหว่างการประมวลผล โมเดลจะส่งคืน Bounding boxes สำหรับวัตถุที่ตรวจจับได้ทั้งหมดจากทั้งสองคลาส ซึ่งจะถูกจับคู่ตามความสัมพันธ์เชิงพื้นที่ (การจับคู่แผ่นยาที่ใกล้ที่สุดกับบริเวณยับยั้ง) เพื่อเชื่อมโยงแต่ละบริเวณยับยั้งเข้ากับแผ่นยาปฏิชีวนะที่เกี่ยวข้อง

\subsection{Hough Circle Transform}

Hough Circle Transform ถูกนำมาใช้เป็นเกณฑ์มาตรฐานแบบคลาสสิกโดยใช้ฟังก์ชัน \texttt{HoughCircles} ของ OpenCV ซึ่งประยุกต์ใช้ Generalized Hough transform สำหรับการตรวจจับวงกลมโดยใช้การลงคะแนนตามความชัน (Gradient-based voting) ก่อนที่จะประยุกต์ใช้การแปลงนี้ รูปภาพแต่ละภาพจะถูกแปลงเป็นสีเทาและทำให้เรียบขึ้นด้วย Gaussian filter โดยใช้ขนาด Kernel $17 \times 17$ พิกเซล และค่า $\sigma = 0$ (เพื่อให้ OpenCV คำนวณค่าเบี่ยงเบนมาตรฐานโดยอัตโนมัติจากขนาด Kernel) เพื่อลดสัญญาณรบกวนที่อาจก่อให้เกิดการตอบสนองต่อความชันที่ผิดพลาด

ฟังก์ชัน \texttt{HoughCircles} รับพารามิเตอร์หลักสี่ตัว ได้แก่ \texttt{dp} (อัตราส่วนความละเอียดผกผันของ Accumulator ต่อภาพนำเข้า ซึ่งควบคุมความละเอียดของ Accumulator ใน Parameter-space), \texttt{minDist} (ระยะห่างขั้นต่ำที่อนุญาตระหว่างจุดศูนย์กลางของวงกลมที่ตรวจจับได้ เพื่อป้องกันการตรวจจับซ้ำซ้อนในบริเวณเดียวกัน), \texttt{param1} (เกณฑ์การตรวจจับขอบ Canny ตัวบน โดยเกณฑ์ตัวล่างจะถูกตั้งเป็นครึ่งหนึ่งของค่านี้โดยอัตโนมัติ) และ \texttt{param2} (เกณฑ์ Accumulator สำหรับการตรวจจับจุดศูนย์กลาง ซึ่งควบคุมจำนวนการลงคะแนนขั้นต่ำที่จำเป็นในการประกาศวงกลมที่ตรวจจับได้) ค่าที่เหมาะสมที่สุดสำหรับพารามิเตอร์ทั้งสี่นี้ได้มาจากการค้นหาแบบตาราง (Grid search) อย่างละเอียดทั่วทั้งชุดข้อมูล เพื่อลดค่าความผิดพลาดสัมบูรณ์เฉลี่ย (Mean absolute error) ระหว่างรัศมีของวงกลมที่ตรวจจับได้กับค่า Ground-truth ค่าสุดท้ายที่เลือกคือ \texttt{dp=1.5}, \texttt{minDist=150}, \texttt{param1=100} และ \texttt{param2=50} สิ่งสำคัญคือตัวเลขประสิทธิภาพสุดท้ายที่รายงานนี้เป็นกรณีที่ดีที่สุดสำหรับ Hough Transform เนื่องจากค่าพารามิเตอร์ถูกปรับจูนโดยใช้ชุดข้อมูลทั้งหมด ซึ่งรวมถึงชุดทดสอบในกระบวนการนี้ด้วย ดังนั้นประสิทธิภาพในการทำงานทั่วไปอาจต่ำกว่าตัวเลขที่รายงาน

\section{การวัดเส้นผ่านศูนย์กลาง ZOI และการปรับเทียบ (ZOI Diameter Measurement and Calibration)}

หลังจากตรวจจับบริเวณยับยั้งในฐานะ Bounding boxes (สำหรับโมเดล Deep Learning ทั้งสอง) หรือเป็นวงกลมที่มีพารามิเตอร์รัศมีที่ชัดเจน (สำหรับ Hough Transform) ผลลัพธ์ดิบจะต้องถูกแปลงเป็นหน่วยวัดเส้นผ่านศูนย์กลาง ZOI ทางกายภาพในหน่วยมิลลิเมตรก่อนที่จะนำไปเปรียบเทียบกับ Breakpoints ของ CLSI หรือ EUCAST

สำหรับการตรวจจับด้วย Bounding box (Faster R-CNN และ YOLOv11n) ระบบประมาณเส้นผ่านศูนย์กลางจาก Bounding box ด้วยสองวิธีดังนี้:
\begin{itemize}
  \item \textbf{อิงตามเส้นผ่านศูนย์กลาง (Diameter-based):} ค่าเฉลี่ยของความกว้าง $w$ และความสูง $h$ ของ Bounding box ในหน่วยพิกเซล คือ $\hat{d}_{\mathrm{diam}} = (w + h)/2$ โดยสมมติว่าบริเวณนั้นเป็นวงกลมโดยประมาณ
  \item \textbf{อิงตามพื้นที่ (Area-based):} เส้นผ่านศูนย์กลางของวงกลมที่มีพื้นที่เท่ากันซึ่งคำนวณจากพื้นที่ Bounding box คือ $\hat{d}_{\mathrm{area}} = 2\sqrt{(w \cdot h)/\pi}$ ซึ่งมีความทนทานมากกว่าสำหรับบริเวณที่มีลักษณะเป็นวงรีเล็กน้อยที่อัตราส่วนภาพของ Bounding box เบี่ยงเบนไปจากหนึ่ง
\end{itemize}

การแปลงจากพิกเซลเป็นมิลลิเมตรดำเนินการโดยใช้แผ่นยาปฏิชีวนะเป็นตัวอ้างอิงในการปรับเทียบภายในภาพ เนื่องจากเส้นผ่านศูนย์กลางทางกายภาพของแผ่นยาปฏิชีวนะมาตรฐานตามข้อกำหนดของ CLSI คือ 6\,มม. พอดี และแผ่นยาก็ถูกตรวจจับเป็น Bounding box โดยโมเดลเช่นกัน ดังนั้นปัจจัยการปรับเทียบ (Calibration factor ในหน่วยมิลลิเมตรต่อพิกเซล) จึงสามารถคำนวณได้สำหรับแต่ละภาพแยกกันดังนี้ $k = 6\,\mathrm{mm} / \hat{d}_{\mathrm{disk}}$ โดยที่ $\hat{d}_{\mathrm{disk}}$ คือค่าประมาณเส้นผ่านศูนย์กลางในหน่วยพิกเซลของ Bounding box ของแผ่นยาปฏิชีวนะที่เกี่ยวข้อง การใช้ปัจจัยการปรับเทียบต่อภาพนี้ช่วยขจัดผลกระทบของระยะห่างของกล้องและระดับการซูมที่แตกต่างกัน ทำให้มั่นใจได้ว่าการวัดมีความแม่นยำแม้จะมีความแปรปรวนของการถ่ายภาพ โดยไม่จำเป็นต้องใช้ไม้บรรทัดจริงในมุมมองของกล้อง

การปรับปรุงการปรับเทียบเป็นกระบวนการที่ทำซ้ำตลอดทั้งโปรเจกต์ การปรับเทียบเบื้องต้นโดยใช้การประมาณเส้นผ่านศูนย์กลางอิงตามพื้นที่เพียงอย่างเดียวส่งผลให้เกิดความคลาดเคลื่อนเชิงระบบในทิศทางบวกประมาณ $+4$\,มม. (ระบบประเมินเส้นผ่านศูนย์กลางบริเวณยับยั้งสูงเกินจริงอย่างต่อเนื่อง) การวิเคราะห์สาเหตุของความคลาดเคลื่อนนี้พบว่า Bounding box ที่พยากรณ์โดยโมเดลมักจะกว้างกว่าขอบเขต ZOI จริงเล็กน้อยเนื่องจากขอบของบริเวณยับยั้งมีความฟุ้งและไล่ระดับสี การผสมผสานการประมาณแบบอิงตามพื้นที่และอิงตามเส้นผ่านศูนย์กลางเข้าด้วยกัน และการประยุกต์ใช้ปัจจัยการแก้ไขเชิงเส้นที่เรียนรู้มา (Learned linear correction factor) ช่วยลดความคลาดเคลื่อนเชิงระบบจากประมาณ $+4$\,มม. เหลือประมาณ $+2.1$\,มม. ซึ่งถือเป็นการปรับปรุงความแม่นยำในการวัดอย่างมีนัยสำคัญ

\section{Pipeline ของระบบ OCR สำหรับระบุชื่อยาปฏิชีวนะ (OCR Pipeline for Antibiotic Label Recognition)}

Pipeline ของระบบ OCR ทำหน้าที่สกัดรหัสตัวอักษรและตัวเลขที่พิมพ์อยู่บนแผ่นดิสก์ยาปฏิชีวนะแต่ละแผ่น จากนั้นนำรหัสที่ได้ไปจับคู่กับรายชื่อยาปฏิชีวนะในฐานข้อมูลเกณฑ์จุดตัด ระบบถูกออกแบบเป็น \textbf{Hybrid Pipeline} ที่ใช้ Gemini 2.0 Flash เป็น OCR Engine หลัก และสลับไปใช้ EasyOCR เป็น Fallback โดยอัตโนมัติเมื่อ API Quota หมดหรือเกิดข้อผิดพลาดในการเชื่อมต่อ

\subsection{การตัดภาพแผ่นดิสก์ยา (Disk Crop Extraction)}

พื้นที่ของแผ่นดิสก์ยาปฏิชีวนะแต่ละแผ่นจะถูกตัด (Crop) แยกออกจากภาพจานเพาะเชื้อเต็มใบ โดยอาศัยพิกัดกรอบ Bounding Box ของคลาส \texttt{antibiotic} ที่โมเดล YOLO ทำนายไว้ พร้อมเผื่อระยะขอบ (Padding) เพื่อให้มั่นใจว่าพื้นที่ผิวแผ่นยาครบถ้วนแม้กรอบจะคลาดเคลื่อนเล็กน้อย ระบบเก็บภาพครอปไว้ \textbf{2 รูปแบบ} คือ ภาพสีต้นฉบับ ($640 \times 640$ พิกเซล) สำหรับส่งไปยัง Gemini API และภาพระดับสีเทาที่ผ่านการประมวลผลเบื้องต้นสำหรับ EasyOCR Fallback

\subsection{ขั้นตอนหลัก: Gemini 2.0 Flash OCR (Primary Stage: Gemini 2.0 Flash OCR)}

ในขั้นตอนแรก ระบบจะส่งภาพสีต้นฉบับของแผ่นดิสก์ในรูปแบบ JPEG ไปยัง Gemini 2.0 Flash API พร้อมกับ Structured Prompt ที่ระบุบริบทอย่างชัดเจนว่า ``แผ่นดิสก์มีรหัสสั้น ประกอบด้วยตัวอักษรพิมพ์ใหญ่แทนตัวย่อยา ตามด้วยตัวเลขแทนขนาดยา (เช่น \texttt{OX 1}, \texttt{AMX 25}, \texttt{P 10}) ให้ตอบกลับเฉพาะรหัสที่พิมพ์บนแผ่นเท่านั้น หากภาพไม่ชัดให้เดาที่ดีที่สุด'' การระบุตัวอย่างรหัสยาจริงใน Prompt ช่วยให้โมเดลจำกัดขอบเขตคำตอบให้ตรงกับรูปแบบที่คาดหวัง และลดโอกาสเกิด Hallucination ได้อย่างมีประสิทธิภาพ

ข้อความที่ได้รับกลับจาก Gemini จะผ่านการตรวจสอบด้วย Regular Expression \texttt{\^{}([A-Za-z]+(?:\textbackslash{}s[A-Za-z]+)\{0,3\})\textbackslash{}s+(\textbackslash{}d+(?:\textbackslash{}.\textbackslash{}d+)?)\$} เพื่อยืนยันว่าผลลัพธ์มีโครงสร้างถูกต้อง (ตัวย่อยาตามด้วยตัวเลข) หากตรงกับ Pattern ระบบจะกำหนดค่าความเชื่อมั่น 0.98 และหากข้อความอยู่นอก Pattern แต่ยังมีเนื้อหาที่สมเหตุสมผล ระบบจะกำหนดค่าความเชื่อมั่น 0.90 เพื่อส่งต่อให้ขั้นตอน Post-processing ตัดสินใจในภายหลัง

เมื่อ Gemini API ส่งกลับ Error หรือตรวจพบสัญญาณ Quota Exhausted (HTTP 429 หรือ \texttt{resource\_exhausted}) ระบบจะตั้งค่า Session Flag เพื่อข้าม Gemini สำหรับทุกแผ่นดิสก์ที่เหลือในเซสชันนั้น และสลับไปใช้ EasyOCR แทนโดยอัตโนมัติ

\subsection{ขั้นตอน Fallback: EasyOCR พร้อมการประมวลผลเบื้องต้น (Fallback Stage: EasyOCR with Preprocessing)}

เมื่อ Gemini API ไม่พร้อมใช้งาน ระบบจะใช้ EasyOCR \cite{jaidedai2024easyocr} ร่วมกับ Pipeline การประมวลผลภาพเบื้องต้น ดังนี้:
\begin{enumerate}
  \item \textbf{CLAHE:} ปรับสมดุลความเปรียบต่างแสงเฉพาะที่ (ขนาด Tile $8 \times 8$ พิกเซล, ค่าจำกัดความเปรียบต่าง 2.0) เพื่อแก้ปัญหาแสงสะท้อนหรือเงาตกกระทบบนพื้นผิวแผ่นดิสก์
  \item \textbf{Otsu Thresholding:} แปลงภาพเป็นขาวดำด้วย Global Thresholding เพื่อดึงตัวอักษรให้เด่นชัดออกจากพื้นหลัง
  \item \textbf{Adaptive Morphological Operations:} หากอัตราส่วนพิกเซลตัวอักษร $< 0.05$ ระบบจะทำ Dilation (Kernel $3 \times 3$) เพื่อเชื่อมเส้นที่ขาดหาย หากอัตราส่วน $> 0.45$ ระบบจะทำ Erosion เพื่อแยกตัวอักษรที่ติดกัน
  \item \textbf{การหมุนภาพ 24 มุม:} EasyOCR รันซ้ำบนภาพที่ผ่านการ Preprocess แล้วครอบคลุม \textbf{24 มุมหมุน} (ตั้งแต่ $-180^\circ$ ถึง $+165^\circ$ เพิ่มทีละ $15^\circ$) เพื่อรับมือกับแผ่นดิสก์ที่ถูกวางในทิศทางใดก็ได้ ผลลัพธ์สุดท้ายจะถูกคัดเลือกจากสตริงที่ผ่านการตรวจสอบ Regex และมีค่าความเชื่อมั่นที่ปรับแก้แล้วสูงที่สุดจากทั้ง 24 มุม
\end{enumerate}

\section{การติดตั้งระบบและโครงสร้างพื้นฐาน (System Deployment and Infrastructure)}

ระบบขั้นสุดท้ายได้รับการเปลี่ยนผ่านจากสภาพแวดล้อมการพัฒนาในเครื่อง (Local) ไปยังการติดตั้งใช้งานในระดับการผลิต (Production-equivalent deployment) ที่เซิร์ฟเวอร์ของ \textbf{ราชวิทยาลัยจุฬาภรณ์ (Chulabhorn Royal Academy หรือ CRA)} โครงสร้างพื้นฐานสำหรับการผลิตนี้ถูกออกแบบมาเพื่อความพร้อมใช้งานสูง (High availability) และการเข้าถึงทางคลินิกที่ปลอดภัย

\subsection{สภาพแวดล้อมเซิร์ฟเวอร์และบริการ Backend (Server Environment and Backend Services)}

บริการ Backend ถูกโฮสต์บนสภาพแวดล้อมเซิร์ฟเวอร์ระบบ Linux ที่ CRA แอปพลิเคชัน FastAPI ให้บริการโดยใช้ \textbf{Uvicorn} ซึ่งเป็น ASGI server ที่รวดเร็วมาก โดยรันเป็นกระบวนการพื้นหลัง (Background process) ที่ได้รับการจัดการ เพื่อให้มั่นใจในเสถียรภาพของระบบ จึงมีสคริปต์จัดการ \texttt{start.sh} โดยเฉพาะที่ทำหน้าที่ประสานงานบริการ, ตรวจสอบกระบวนการ และกู้คืนระบบโดยอัตโนมัติ Backend ถูกกำหนดค่าให้ฟังคำขอ (Listen) บนพอร์ตหมายเลขสูง (6014) เพื่อหลีกเลี่ยงความขัดแย้งกับบริการอื่นๆ ของโรงพยาบาล

\subsection{Reverse Proxy และความปลอดภัย (Reverse Proxy and Security)}

มีการติดตั้ง \textbf{Nginx} เป็น Reverse proxy และ Web server เพื่อจัดการการรับส่งข้อมูลที่เข้ามายังแอปพลิเคชัน Nginx ทำหน้าที่จัดการ SSL/TLS termination เพื่อให้สามารถเข้าถึงระบบผ่าน HTTPS ได้อย่างปลอดภัย โดยมีการกำหนดค่าบล็อก \texttt{location} เฉพาะสำหรับการกำหนดเส้นทาง:
\begin{itemize}
    \item \texttt{/zoneanalyzer/}: ส่งต่อข้อมูลไปยัง React frontend (npx serve บนพอร์ต 6015)
    \item \texttt{/zoneanalyzer/api/}: ส่งต่อคำขอไปยัง FastAPI backend (พอร์ต 6014)
    \item \texttt{/zoneanalyzer/uploaded\_images/}: ให้บริการรูปภาพที่แสดงผลลัพธ์โดยตรงจากพื้นที่จัดเก็บของเซิร์ฟเวอร์
\end{itemize}
การกำหนดค่านี้ช่วยให้ใช้จุดเชื่อมต่อโดเมนเดียว (\texttt{https://paniti-jupyter.cra.ac.th/zoneanalyzer/}) สำหรับทั้งแอปพลิเคชัน ช่วยลดความยุ่งยากในการติดตั้งทางคลินิกและการจัดการไฟร์วอลล์

\subsection{การให้บริการ Frontend (Frontend Serving)}

React frontend ได้รับการปรับแต่งเพื่อการผลิตโดยใช้ Vite build pipeline ซึ่งจะสร้างชุดไฟล์ Static assets ที่มีขนาดเล็กและผ่านการทำ Tree-shaken ไฟล์เหล่านี้ถูกให้บริการโดยใช้ \textbf{npx serve} ซึ่งเป็นกลไกการส่งมอบอินเทอร์เฟซผู้ใช้ที่มีน้ำหนักเบาและมีประสิทธิภาพ สภาพแวดล้อม Frontend ถูกกำหนดค่าด้วยไฟล์ \texttt{.env} เฉพาะที่ชี้ไปยัง Endpoint ของ Production API เพื่อให้แน่ใจว่ามีการสื่อสารที่ราบรื่นผ่าน Reverse proxy

\subsection{การประมวลผลข้อความและจับคู่ชื่อยาปฏิชีวนะ (Post-Processing and Antibiotic Name Matching)}

ผลลัพธ์ข้อความดิบที่ได้จาก OCR จะถูกกรองผ่านนิพจน์ปกติ (Regular Expression หรือ Regex) เพื่อปฏิเสธสตริงที่มีรูปแบบไม่ถูกต้อง โดยใช้แพตเทิร์น \texttt{\^{}([A-Za-z]+(\textbackslash{}s[A-Za-z]+)\{0,3\})\textbackslash{}s+(\textbackslash{}d+)\$} เพื่อบังคับโครงสร้างข้อมูลให้ประกอบด้วยตัวอักษร 1--4 ตัว (แทนตัวย่อของยา) ตามด้วยตัวเลข (แทนความเข้มข้นในหน่วยไมโครกรัม) สตริงที่ไม่ตรงกับรูปแบบนี้ ไม่ว่าจะเป็นข้อความที่ระบบหลอนขึ้นมาเอง (Hallucinations), การอ่านค่าได้เพียงบางส่วน หรือสตริงที่มีแต่ตัวเลข จะถูกระบบตัดทิ้ง และแผ่นยาตรงจุดนั้นจะถูกระบุสถานะในผลลัพธ์ว่า ``ไม่สามารถระบุได้ (Unidentified)''

รหัสตัวย่อยาที่ผ่านการกรอง (เช่น ``OX'', ``AMX'', ``CIP'') จะถูกนำไปค้นหาในตารางจับคู่ชื่อยาปฏิชีวนะ เพื่อแปลงเป็นชื่อเต็มและดึงข้อมูลประเภทของยา กระบวนการจับคู่นี้เป็นขั้นตอนสำคัญที่ช่วยให้ระบบสามารถไปคิวรีเกณฑ์จุดตัด (Breakpoints) จากฐานข้อมูล CLSI หรือ EUCAST ที่ถูกต้องตรงตามชนิดของยาและกลุ่มเชื้อแบคทีเรียที่ผู้ใช้กำหนดได้

\section{การจำแนกความไวต่อยา (Susceptibility Classification)}

เพื่อให้ระบบสามารถทำหน้าที่แปลผลได้อย่างสมบูรณ์ จึงได้นำตรรกะการตัดสินผล (Breakpoint Logic) อ้างอิงจากมาตรฐาน CLSI 2025 \cite{clsi2025} และ EUCAST 2023 \cite{eucast2023} เข้ามาบูรณาการในระบบ สำหรับบริบทของห้องปฏิบัติการคลินิกในประเทศไทย มาตรฐานหลักที่ดูแลโดยกรมวิทยาศาสตร์การแพทย์ สถาบันวิจัยวิทยาศาสตร์สาธารณสุข (NIH Thailand) \cite{dmsc_nih} มีแนวทางการประเมินผลที่สอดคล้องอย่างใกล้ชิดกับแนวทางของ CLSI โครงงานนี้จึงเลือกใช้ CLSI เป็นมาตรฐานอ้างอิงหลัก ฐานข้อมูลเกณฑ์จุดตัดถูกจัดเก็บไว้ในตาราง SQLite โดยมีโครงสร้าง (Schema) ประกอบด้วย: ชื่อยาปฏิชีวนะ, กลุ่มเชื้อแบคทีเรีย, มาตรฐานที่ใช้อ้างอิง, ความเข้มข้นของแผ่นยา (Potency), ขนาดเส้นผ่านศูนย์กลางเกณฑ์ไว ($S_{\min}$), ช่วงกึ่งไว ($I_{\min}$--$I_{\max}$ ถ้ามี) และขนาดเส้นผ่านศูนย์กลางเกณฑ์ดื้อ ($R_{\max}$) ซึ่งโครงสร้างนี้สะท้อนข้อมูลจากตารางมาตรฐานฉบับพิมพ์โดยตรง

\begin{figure}[H]
  \centering
  \safeimg[width=\textwidth]{figures/_review/database_schema.png}{Database ER diagram}
  \caption{แผนภาพ Entity-Relationship (ER) ของฐานข้อมูล SQLite แสดงความสัมพันธ์ระหว่างตาราง
    Users, AnalysisBatch, Plates, PlateResults, Breakpoints\_DiskDiffusion, Microbes,
    Standards และ Antibiotics}
  \label{fig:db_schema}
\end{figure}

สำหรับคู่แผ่นยาและ ZOI แต่ละคู่ที่ตรวจจับได้ ระบบจะดำเนินการจำแนกผลตามขั้นตอนดังนี้:
\begin{enumerate}
  \item ระบบคิวรีข้อมูลเกณฑ์จุดตัดที่ตรงกับรหัสยา (จาก OCR), กลุ่มเชื้อที่ระบุ และมาตรฐานที่เลือก (CLSI หรือ EUCAST)
  \item นำขนาดเส้นผ่านศูนย์กลาง ZOI ที่ผ่านการปรับเทียบแล้ว ($d_{\mathrm{mm}}$) ไปเปรียบเทียบกับเกณฑ์ที่ดึงมาจากฐานข้อมูล
  \item ตัดสินผลการจำแนก: เป็น ไวต่อยา (S) หาก $d_{\mathrm{mm}} \geq S_{\min}$, ดื้อยา (R) หาก $d_{\mathrm{mm}} \leq R_{\max}$, หรือ กึ่งไว (I) หากค่าที่วัดได้ตกอยู่ในช่วง Intermediate ของคู่นั้นๆ
  \item ในกรณีที่ไม่พบข้อมูลเกณฑ์จุดตัดที่ตรงกับคู่ยาและเชื้อที่ระบุ ระบบจะส่งคืนค่าเป็น ``ไม่พบข้อมูล (Not Available หรือ NA)'' ทว่ายังคงรายงานค่าเส้นผ่านศูนย์กลางดิบเพื่อให้ผู้ปฏิบัติงานสามารถนำไปเทียบกับคู่มือด้วยตนเองได้
\end{enumerate}

ผลลัพธ์การวิเคราะห์ทั้งหมดบนจานเพาะเชื้อจะถูกรวบรวมและส่งกลับไปยังฝั่ง Frontend ในรูปแบบ JSON โดยหน้าแดชบอร์ด React จะนำไปแสดงผลเป็นรายงานตารางที่ระบุชื่อยา, ความเข้มข้น, ขนาด ZOI และระดับความไวต่อยา (S/I/R) นอกจากนี้ แดชบอร์ดยังแสดงภาพถ่ายจานเพาะเชื้อพร้อมวาดกรอบ Bounding box และภาพซ้อนทับ (Overlay) เพื่อให้ผู้ใช้สามารถตรวจสอบความถูกต้องเชิงพื้นที่ได้ทันที ข้อมูลทั้งหมดจะถูกบันทึกลงในฐานข้อมูล SQLite พร้อมกันเพื่อประโยชน์ในการสืบค้นย้อนหลังและการจัดการข้อมูลในห้องปฏิบัติการอย่างมีประสิทธิภาพ
